% Chapter 1: Probability and inference
% Bayesian Data Analysis - Gelman et al.

\chapter{概率与推断}

\section{贝叶斯数据分析的三步骤}

贝叶斯数据分析的过程可以理想化地分为以下三个步骤：

\begin{enumerate}
  \item \textbf{建立完整的概率模型}：对问题中所有可观测和不可观测的量建立一个联合概率分布。模型应与底层科学问题（scientific problem）和数据收集过程的知识相一致。

  \item \textbf{基于观测数据进行分析}：计算并解释适当的后验分布——即在给定观测数据的条件下，关于不可观测量的条件概率分布。

  \item \textbf{评估模型的拟合度和后验分布的含义}：
  \begin{itemize}
    \item 模型对数据的拟合程度如何？
    \item 实质性结论（substantive conclusions）是否合理？
    \item 结果对第一步中建模假设的敏感性如何？
  \end{itemize}
  作为回应，可以修改或扩展模型并重复这三个步骤。
\end{enumerate}

第二个步骤涉及计算方法论，第三个步骤涉及微妙的技术与判断平衡，均由问题的应用背景指导。

\textbf{第一步仍然是许多贝叶斯分析的主要障碍}：我们的模型从何而来？如何构建适当的概率规范？

贝叶斯思维的主要动机是它便于对统计结论进行常识性解释。例如，一个未知量的贝叶斯（概率）区间可以直接被认为是"该未知量有很高概率包含在其中"，这与频率主义（置信区间）的解释形成对比——后者在严格解释下只能关联到重复抽样序列中的类似推断。

\section{统计推断的一般符号}

在贝叶斯统计中，我们使用以下符号：

\begin{itemize}
  \item $\theta$：不可观测的向量值或总体参数（如临床试验中每种治疗方案的生存概率）

  \item $y$：观测数据（如每组中的存活和死亡人数）

  \item $\tilde{y}$：未知但潜在可观测的量（如其他治疗方案下患者的结果，或与试验中患者相似的新患者的结果）

  \item $x$：解释变量或协变量（如患者的年龄和先前健康状况）
\end{itemize}

这些概率陈述以观测值 $y$ 为条件，在我们记号中简单记为 $p(\theta|y)$ 或 $p(\tilde{y}|y)$。我们也隐含地以协变量的已知值为条件。

\section{贝叶斯推断}

贝叶斯推断的核心思想是：关于参数 $\theta$ 或未观测数据 $\tilde{y}$ 的统计结论，都以概率陈述的形式给出。

这些概率陈述以观测值 $y$ 为条件，在我们的记号中简单记为 $p(\theta|y)$ 或 $p(\tilde{y}|y)$。

\subsection{贝叶斯定理}

为了对给定 $y$ 的 $\theta$ 做出概率陈述，我们必须首先建立一个提供 $\theta$ 和 $y$ 联合概率分布的模型。

联合概率质量（或密度）函数可以写成两个密度的乘积，通常称为先验分布 $p(\theta)$ 和采样分布（或数据分布）$p(y|\theta)$：
\[
p(\theta, y) = p(\theta) \times p(y|\theta)
\]

使用条件概率的基本性质（贝叶斯定理），得到后验密度：\footnote{推导见附录 \cref{sec:bayes-theorem-derivation}}
\[
p(\theta|y) = \frac{p(\theta) \times p(y|\theta)}{p(y)}
\]
其中
\[
p(y) = \int p(\theta) \times p(y|\theta) \, d\theta
\]

一个等价的形式省略了不依赖于 $\theta$ 的因子 $p(y)$，给出非归一化后验密度：
\[
p(\theta|y) \propto p(\theta) \times p(y|\theta)
\]

\textbf{为什么可以省略 $p(y)$？}

关键在于：后验分布的\textbf{核}（kernel）\footnote{关于"核"概念的详细解释，见附录 \cref{sec:qa:DistributionKernel}}包含了所有关于 $\theta$ 的信息，而 $p(y)$ 作为归一化常数不包含任何关于 $\theta$ 的新信息。

以正态分布为例：后验密度 $p(\theta|y) \propto \exp\left(-\frac{(\theta-\mu)^2}{2\sigma^2}\right)$，我们只凭这个指数函数的形式就能识别出这是正态分布 $N(\mu, \sigma^2)$——因为已知核是二次型的指数，就能确定完整分布。

换句话说，$p(y)$ 的作用只是确保后验密度积分为 1，但它的具体值对于理解 $\theta$ 如何随数据更新并无贡献。

\begin{Remark}
  注意：$p(y|\theta)$ 被视为 $\theta$ 的函数（固定 $y$），而不是 $y$ 的函数。这被称为似然函数。
\end{Remark}

\subsection{预测}

在数据 $y$ 被观测之前，未知但可观测的 $\tilde{y}$ 的分布是边际分布：\footnote{推导见附录 \cref{sec:predictive-distributions-derivation}}
\[
p(y) = \int p(y|\theta) \times p(\theta) \, d\theta
\]
这被称为先验预测分布（prior predictive distribution）。

在观测到数据 $y$ 之后，预测 $\tilde{y}$ 的分布是后验预测分布（posterior predictive distribution）：\footnote{推导见附录 \cref{sec:predictive-distributions-derivation}}
\[
p(\tilde{y}|y) = \int p(\tilde{y}|\theta) \times p(\theta|y) \, d\theta
\]

后验预测分布可以解释为在 $\theta$ 的后验分布上对条件预测进行平均。

\begin{Example}[称重例子]
  令 $y = (y_1, \ldots, y_n)$ 为物体在秤上称量 $n$ 次的记录重量，$\theta = (\mu, \sigma^2)$ 为物体的真实重量和测量方差，$\tilde{y}$ 为计划中新的称量结果。

  第二行展示了后验预测分布作为后验分布上条件预测的平均。最后一步来自在给定 $\theta$ 的条件下 $y$ 和 $\tilde{y}$ 的条件独立性假设。
\end{Example}

\section{似然与优势比}

使用贝叶斯定理与选定的概率模型意味着数据 $y$ 只通过 $p(y|\theta)$ 影响后验推断，当固定 $y$ 将其视为 $\theta$ 的函数时，称为\textbf{似然函数}。

这样，贝叶斯推断服从所谓的似然原理：对于给定的数据样本，任何具有相同似然函数的概率模型 $p(y|\theta)$ 对 $\theta$ 产生相同的推断。

后验密度 $p(\theta|y)$ 在 $\theta_1$ 和 $\theta_2$ 处取值的比率称为后验优势比。优势比提供了一个概率的替代表示，并且具有贝叶斯定理用优势比表示特别简单的形式：\footnote{推导见附录 \cref{sec:posterior-odds-derivation}}
\[
\frac{p(\theta_1|y)}{p(\theta_2|y)} = \frac{p(\theta_1)}{p(\theta_2)} \times \frac{p(y|\theta_1)}{p(y|\theta_2)}
\]

即：\textbf{后验优势比 = 先验优势比 × 似然比}。

\section{概率作为不确定性的度量}

在贝叶斯统计中，概率被用作不确定性的基本度量（fundamental measure or yardstick of uncertainty）。在这个范式下，讨论明天降雨的概率或巴西队赢得世界杯的概率与讨论硬币抛掷正面朝上的概率同样合法。

\textbf{为什么概率是量化不确定性的合理方式？}通常有以下理由：

\begin{enumerate}
  \item \textbf{类比}：物理随机性引起不确定性，因此用随机事件的语言描述不确定性似乎是合理的。日常语言使用"可能"和"不太可能"等术语，将更正式的概率计算扩展到科学推断问题是与之一致的。

  \item \textbf{公理化或规范性方法}：与决策理论相关，将所有统计推断置于收益和损失的决策背景下。那么合理的公理（排序、传递性等）意味着不确定性必须用概率表示。

  \item \textbf{博彩一致性}：定义你赋予事件 $E$ 的概率 $p$（$p \in [0, 1]$）为你愿意以此比例交换（即赌）\$p 来换取如果 $E$ 发生则获得\$1 的回报。也就是说，如果 $E$ 发生，你额外获得 $(1-p)$；如果 $E$ 的补集发生，你损失 $p$。
\end{enumerate}

一致性原则（coherence principle）表明：你对所有可能事件的概率赋值不应该让对方有可能从你这里确定获利。可以证明在这种原则下构建的概率必须满足概率论的基本公理。

博彩基本原理有一些根本性困难：
\begin{itemize}
  \item 需要精确的赔率，你愿意在两个方向中的任何一个方向下赌注。对于所有事件如何指定精确赔率？如果你不确定呢？
  \item 如果一个人愿意与你对赌，而且他有你不知道的信息，你接受赌注可能不明智。
\end{itemize}

这些考虑表明概率可能是总结应用统计中不确定性的合理方法，但最终证明在于应用的成功。

\subsection{两种经典论证}

\textbf{对称性/可交换性论证}：
\[
\text{概率} = \frac{\text{有利情况数}}{\text{可能情况数}}
\]

对于硬币抛掷，这实际上是一个物理论证，基于关于决定硬币如何落地的力量以及抛掷的初始物理条件的假设。

\textbf{频率论证}：
\[
\text{概率} = \text{在长期重复序列中得到的相对频率}
\]

假设以相同方式物理独立地进行。

这两种论证在某种意义上都是主观的，因为它们需要对硬币性质和抛掷程序的判断，两者都涉及关于等可能性事件、相同测量和独立性的语义论证。

频率论证可能有某些特殊困难，因为它涉及长期相同抛掷序列的假设性概念。如果严格遵循，这种观点不允许对不属于长期相同事件序列的单一抛掷做出概率陈述。

\section{主观性与客观性}

所有使用概率的统计方法在某种意义上都是主观的，因为它们依赖于对世界的数学理想化。

贝叶斯方法有时被认为特别主观，因为它们依赖于先验分布。但在大多数问题中，科学判断对于指定似然和先验分布都是必要的。

一个普遍原则是：每当存在复制时（即观察到许多可交换的单位），就有机会从数据中估计概率分布的特征，从而使分析更加客观。

如果整个实验被重复若干次，那么先验分布的参数本身也可以从数据中估计。

\section{概率赋值的例子：足球点差}

在足球比赛中，"点差"（point spread）是博彩公司设定的赔率，用于使两边下注相对均衡。例如，如果热门队（favorite）的点差是 8 分，意味着博彩公司认为热门队平均会赢 8 分。

给定点差，我们可以估计各种条件概率：
\begin{itemize}
  \item $\Pr(\text{热门队赢} | \text{点差}=8)$
  \item $\Pr(\text{热门队至少赢 8 分} | \text{点差}=8)$
  \item $\Pr(\text{热门队至少赢 8 分} | \text{点差}=8 \text{ 且 热门队赢})
\end{itemize}

这些概率可以通过相对频率估计，也可以用正态近似来估计（outcome - point spread 的分布）。

\section{概率论的一些有用结果}

见教材。

\section{本章小结}

本章我们学习了：
\begin{enumerate}
  \item 贝叶斯数据分析的三个基本步骤
  \item 统计推断的一般符号 ($\theta, y, \tilde{y}, x$)
  \item 贝叶斯定理及其组成部分（先验、似然、后验、边际似然）
  \item 预测分布：先验预测分布与后验预测分布
  \item 概率作为不确定性度量的多种解释
  \item 似然原理和优势比
\end{enumerate}

\section{下节预告}

下一章我们将学习单参数模型，包括从二项数据估计概率的经典例子——这是贝叶斯推断最经典的入门案例。

\endinput
